% BHU PYQ -- Manually transcribed & error-corrected
\documentclass[11pt,a4paper]{article}

\usepackage[a4paper,top=2.5cm,bottom=2.5cm,left=2.8cm,right=2.8cm,headheight=14pt]{geometry}
\usepackage{amsmath,amssymb}
\usepackage[T1]{fontenc}
\usepackage[utf8]{inputenc}
\usepackage{lmodern}
\usepackage{microtype}
\usepackage{fancyhdr}
\usepackage{enumitem}
\usepackage{hyperref}
\usepackage{lastpage}
\usepackage{parskip}

\hypersetup{colorlinks=true,urlcolor=black,linkcolor=black,
  pdftitle={BPT-604: Atomic Physics and Lasers -- BHU PYQ},pdfauthor={Banaras Hindu University}}

\pagestyle{fancy}\fancyhf{}
\renewcommand{\headrulewidth}{0.4pt}\renewcommand{\footrulewidth}{0.4pt}
\fancyhead[L]{\small\textit{Banaras Hindu University}}
\fancyhead[R]{\small\textit{BPT-604: Atomic Physics and Lasers}}
\fancyfoot[C]{\small Page \thepage\ of \pageref{LastPage}}

\newlist{parts}{enumerate}{2}
\setlist[parts,1]{label=(\alph*),leftmargin=*,itemsep=5pt,topsep=3pt}
\setlist[parts,2]{label=(\roman*),leftmargin=*,itemsep=3pt,topsep=2pt}
\newcommand{\pts}[1]{\hfill\textit{[#1]}}

\begin{document}

\begin{center}
  {\large\bfseries BANARAS HINDU UNIVERSITY}\\[2pt]
  {\normalsize B.Sc. (Hons.) Semester VI Examination, 2013-14}\\[6pt]
  \rule{0.8\linewidth}{0.5pt}\\[6pt]
  {\large\bfseries Paper No. BPT-604: Atomic Physics and Lasers}\\[4pt]
  {\normalsize Time: 3 Hours\hfill Maximum Marks: 70}
\end{center}
\rule{\linewidth}{0.4pt}
\smallskip
\noindent\textit{Instructions: Answer all questions. Symbols have their usual meanings.}
\bigskip

\noindent\textbf{Question 1.} Answer any \textbf{five} of the following: \pts{5 $\times$ 4 = 20}
\begin{parts}
  \item What is the relative population of the two states in a ruby laser that produces a light beam of wavelength $6943\,\text{\AA}$ at $27^\circ\text{C}$?
  \item Why is the discharge tube in a He-Ne laser system taken to have a narrow cross-section?
  \item The effect of nuclear mass on the Rydberg constant is to change it by about 1 part in 2000. Justify this statement.
  \item Which experiment provided evidence for the existence of electron spin, and how?
  \item Explain the meaning of space quantization.
  \item On introducing the concept of the relativistic electron in the elliptic Sommerfeld orbit treatment, the splitting of energy levels is the same as obtained by Dirac's theory. Do these two explain the observed structure completely? If not, why?
  \item Work out the ground state of an atom having atomic number 39.
\end{parts}

\medskip\rule{\linewidth}{0.2pt}

\noindent\textbf{Question 2.}
\begin{parts}
  \item Describe the construction and working of an Nd-YAG laser. \pts{12.5}
\end{parts}
\hfill \textbf{OR}
\begin{parts}
  \item With the help of a suitable diagram, describe the construction and working of a He-Ne laser. Which one of the He and Ne atoms is responsible for the LASER action and what is the role of the other? \pts{12.5}
\end{parts}

\medskip\rule{\linewidth}{0.2pt}

\noindent\textbf{Question 3.}
\begin{parts}
  \item What is the difference between spontaneous and stimulated emissions? Establish the relation between the Einstein coefficients $A$ and $B$. \pts{12.5}
\end{parts}
\hfill \textbf{OR}
\begin{parts}
  \item An emitting atom is placed in a magnetic field. If the field intensity is increased from zero to a low value (a few Gauss), the singlet lines split into three components. Explain the origin of this triplet. \pts{12.5}
\end{parts}

\medskip\rule{\linewidth}{0.2pt}

\noindent\textbf{Question 4.}
\begin{parts}
  \item What do you understand by Land\'{e}’s interval rule and what are normal and inverted multiplets? Explain with suitable examples. \pts{7.5}
  \item Find the ground state of an atom with two equivalent $p$ electrons. \pts{5}
\end{parts}

\medskip\rule{\linewidth}{0.2pt}

\noindent\textbf{Question 5.}
\begin{parts}
  \item Describe the Franck-Hertz experiment and its outcome. \pts{4.5}
  \item Discuss briefly the condition of phase matching in Second Harmonic Generation (SHG). \pts{4}
  \item Derive terms for $5d6p$ electrons under $j-j$ coupling scheme. \pts{4}
\end{parts}

\vfill
\rule{\linewidth}{0.4pt}
\begin{center}\small
  \href{https://bhu-exam.vercel.app/}{Click here to visit our website}\\[4pt]
  \href{https://me-aryan.vercel.app/}{Click here to visit our contributor}\\[4pt]
  \href{https://quashan-me.vercel.app/}{Click here to visit our contributor}
\end{center}

\end{document}
